\Needspace{5\baselineskip}\pdfbookmark[2]{First division}{witwen-II-1}
\begin{center}\textit{First division.}\end{center}
\noindent\hspace*{\parindent}The Senate's discussion of the condition of the widows' fund, and the deliberations which led the University Board of Curators to order a thorough investigation of that condition, were initially prompted by the concerns that arose in the autumn of last year. These were aroused especially by the rapid rise in the number of widows, far exceeding all previous experience (to twenty-two since the death of Privy Councillor of Justice M.), and were further heightened by a momentary shortage of cash in the fund to pay the pensions in full on the customary day.

That these and other circumstances caused some uneasiness is all the less surprising since people had grown accustomed to regarding the fund's condition as highly flourishing. Until Easter 1829 the annual pension had been 210 Thl., and through four increases of 10 Thl. each within the short period of \(6\tfrac12\) years by Michaelmas 1835 it had risen to 250 Thl. At that time, therefore, the point when the pensions would have grown to 300 Thl. was thought so near that plans were already being considered for the best use of the surplus thereafter\footnote{Two well-meaning members of the Church Deputation, both since deceased, suggested, the one abolishing the annual contributions, the other extending the orphans' pensions as far as establishing lifelong pensions for the unmarried daughters of professors.}. Yet it was precisely from that time onwards that circumstances began to become less favourable; to the twelve pensioners present at the end of 1835 twelve new widows were added within 9 years, while only two pensions ceased.

The momentary insufficiency of cash in the fund mentioned above is, incidentally, a matter of little importance, even if it had made a short-term interest-bearing loan necessary; according to the annual accounts for 1844—1845, however, this does not seem to have been the case on that occasion. Such eventualities can occur in the soundest fund, as with the soundest private individual engaged in many financial transactions, and all the more often the more effort is made to avoid leaving large sums of money unused for long.

Even the greatly increased number of widows in 1844 was, considered in itself, still no proof of imminent danger. It could reasonably be expected that this number would soon fall again, as indeed three widows died between March and June of this year. Nor is it unlikely that some further decline may occur before long: nevertheless, a calculation from today to tomorrow offers only very inadequate reassurance, and there is certainly no reason to expect a substantial and lasting fall in the number of widows.

It is not the temporary size of the number of widows that threatens the institution with danger, but something quite different, namely
\SourceLeaf{132}
\begin{quote}\textit{the unclear wording of that part of the regulations which is intended to determine the progression of the pension rates}\end{quote}
in conjunction with
\begin{quote}\raggedright\textit{the present greatly increased number of professors participating in the widows' fund.}\end{quote}
\noindent\hspace*{\parindent}The great importance of the latter circumstance has already been indicated in my opinion mentioned above. The significance of a large number of participants differs greatly according to the nature of a widows' fund. For a widows' fund which supports itself entirely through the members' contributions (or through entrance fees, or through both together), a very large number of participants will only be advantageous, provided that the fund rests on a correct calculation. A society of \textit{this} kind with twice as many members must, other circumstances being equal, expect twice as many widows as one of the original size: but it also has exactly twice the income, and can therefore grant each widow exactly as much, with greater protection, however, against the changing fluctuations, which are smaller in relation to the whole in the larger society than in the smaller one.

The situation is quite different, however, with an institution for the support of widows which is a pure benefaction and whose resources have a given size once and for all (through a specified holding of capital or land). Here too, every extension of participation will result in a proportionate increase in the number of widows, so that each individual widow will no longer be able to receive as much from the given resources as she could before, when participation was limited. Admittedly, the increase in the number of widows corresponding to the increased number of participants will take effect to its \textit{full} extent only after several decades and will, in the natural course of things, develop gradually until then. To make the idea definite, let us suppose that the size of such a society (which is a pure benefaction) has doubled within a certain time. One must then be prepared for an increased number of widows before long, and thus, if the assets themselves are not to be drawn upon, for a reduction in the pension to be granted to each individual widow; this reduction will gradually proceed as far as one half. But if such a society had a statute under which pensions of an unchanged amount \textit{must} continue to be paid to the widows despite their growing number, it would necessarily be ruined. There are, however, two cases in which this course of events will or may undergo a modification. \textit{First}, if the fund's resources before participation was extended were more than sufficient to meet the existing pensions, so that continual growth of the assets, and perhaps also an increase in the pension rate from time to time up to that point, had been possible. Here the effect of the extension will clearly depend on \textit{how much?} If the fund previously had a large annual surplus and the increase in the number of participants is not very considerable, it may perhaps survive the danger; the growth of its assets will merely become slower and slower, and it is possible that, when the effect of that cause has developed fully, the fund may still have sufficient means to grant the full pension even to the larger number of widows. Conversely, however, if the annual surplus was initially small but the increase in the number of participants very considerable, the annual surplus will soon turn into a deficit, and the fund's ruin will occur somewhat later than if resources and claims had originally been in equilibrium, but just as inevitably. \textit{Second}, in a fund of small size overall in terms of the number of participants, where this number must therefore still be regarded as small even after the increase, one cannot expect as regular a course of events as in larger funds; the fluctuations inherent in the nature of the matter, but \textit{relatively} much larger with small numbers, may weaken the regularity in the succession of events,
\SourceLeaf{133}
indeed obscure it entirely, without in the least impairing the correctness of the proposition that, according to averages from sufficiently long periods, twice the number of participants must also be followed by twice the number of widows. But for that reason the proportionately increased number of widows may take longer to appear in a small society than in a large one; it may just as well, however, appear much earlier. In a small society it may happen that, for a considerable number of years after the increase in the number of participants, the number of widows shows only a quite insignificant increase, remains almost stationary, or even falls somewhat again. This, however, would in fact be anything but desirable if it lulled the administration into a false sense of security and, relying on the assets still increasing at that moment, it ordered a further \textit{increase} in the pension at a time when a thorough consideration looking beyond the next day might already have recognised the necessity of a \textit{restriction}. For it needs no proof that restrictions which become necessary must be all the greater the longer they have been deferred.

What has been said about funds providing purely gratuitous support can now readily be applied to those which stand between them and funds supporting themselves entirely through contributions. The professors' widows' fund is such a mixed fund, although, because the contributions are so small, it stands much closer to the former than to the latter. On the basis of annual contributions of 10 Thl., as follows from the calculations to be discussed in the second division, the participants' survivors could be granted a pension of at most 44 Thl. or 48 Thl., according as an interest rate of \(3\tfrac12\) or 4 percent is assumed; and nothing has yet been deducted here for possible losses or for administrative and other expenses. Anything granted beyond this, thus, at the pension rate in force since 1835, an annual 202 to 206 Thl., must be regarded as the proceeds of a pure benefaction; with regard to the effects of the growing number of participants, exactly the same applies to it as was developed above concerning such funds.

This does indeed give grave importance to the circumstance that the number of participants in our widows' fund is now greater by one half than it was on average 20 to 30 years ago. To assess this importance properly, however, it must also be carefully considered that the recently so large number of widows, or more correctly of pensions (according to the \textit{average} of the last \textit{eight} years \(=20\)), is in no way a consequence of the present large number of participants, but would be just as large even if the number of participants had not increased so much. This is evident from the fact that the husbands of the widows who have made up the roll over the last eight years (or the fathers of the orphans who received pensions) had almost all belonged to the society before the rise in the number of participants, indeed mostly long before that rise. Rather, in a society whose size is increasing, the number of widows at any given time must be compared with the number of participants several decades earlier, rather than with their number at that same time. The following inference is then very natural: just as the society's former condition, with a number of participants which 20 to 30 years ago fluctuated between 31 and 38, has now produced an average number of 20 widows, so its present size — 51 participants — may quite reasonably produce 30 widows after a further few decades, with no guarantee whatever that this number is a maximum which cannot be exceeded. This does not say that it certainly \textit{will} actually happen, but only that, according to previous precedents, it may happen without our being entitled to regard it as anything extraordinary; in any case, even such a rough estimate shows that, amid the possible changes of fortune, the widows' fund is playing for much higher stakes than may previously have been believed.
